\documentclass[../main.tex]{subfiles}

\begin{document}
\subsection{Funciones armónicas}
\classdate{2026-04-10}

\subsection{Teorema Fundamental del Cálculo Complejo}

\begin{proposition}
    Si \(f \colon \Omega \subseteq \mathbb{C} \to \mathbb{C}\) es continua, \(\gamma
    \colon  \bigl[a, b\bigr] \subseteq \mathbb{R} \to \Omega\) es suave por pedazos
    \(\exists F \colon \Omega \subseteq \mathbb{C} \to \mathbb{C}\) analítica tal que
    \(F^{\prime}(z) = f(z)\; \forall z\in\Omega\),

    \begin{equation*}
        \int_{\gamma}^{}f(z)\odif{z} = \int_{\gamma}^{}F^{\prime}(z)\odif{z} = F \bigl(\gamma(b)\bigr) - F \bigl(\gamma(a)\bigr).
    \end{equation*}
\end{proposition}

\begin{proof}
    \begin{align*}
        \int_{\gamma}^{}f(z)\odif{z} & = \int_{a}^{b}f
        \bigl(\gamma(t)\bigr)\gamma^{\prime}(t)\odif{t} =
        \int_{a}^{b}F^{\prime}\bigl(\gamma(t)\bigr)\gamma^{\prime}(t)\odif{t} =
        \int_{a}^{b}\bigl(F \circ \gamma\bigr)^{\prime}(t)\odif{t},                                         \\
        & = \int_{a}^{b}\Bigl[\Re \bigl(\bigl(F \circ \gamma\bigr)^{\prime}(t)
            \bigr) + i\Im \bigl(\bigl(F \circ \gamma\bigr)^{\prime}(t)
        \bigr)\Bigr]\odif{t},                                                                               \\
        & = \int_{a}^{b}\Re \bigl(\bigl(F \circ \gamma\bigr)^{\prime}(t)
        \bigr)\odif{t} + i\int_{a}^{b}\Im \bigl(\bigl(F \circ \gamma\bigr)^{\prime}(t)
        \bigr)\odif{t},                                                                                     \\
        & = \Re \bigl(F \bigl(\gamma(b)\bigr)\bigr) - \Re \bigl(F
        \bigl(\gamma(a)\bigr)\bigr) + i \Bigl[\Im\bigl(F\bigl(\gamma(b)\bigr)\bigr) -
        \Im\bigl(F \bigl(\gamma(a)\bigr)\bigr)\Bigr],                                                       \\
        & = F(\gamma(b)) - F(\gamma(a)).
    \end{align*}
\end{proof}

\begin{example}
    Sea \(f(z) = z^{2}\) y \(\gamma\) la curva que empieza en \(\pi\) y recorre la
    circunferencia de radio \(\pi\) en sentido antihorario hasta \(-\pi\) y de ahí
    recorre el eje real hasta \(\pi\). Calcule \(\int_{\gamma}^{}f(z)\odif{z}\).
\end{example}

\startsolution
Primero definamos la curva, la cual será una función por partes,

\begin{equation}
    \gamma(t) =
    \begin{dcases}
        \pi \mathrm{e}^{it}, & t\in [0,\pi],    \\
        t - 2\pi,            & t\in[\pi, 3\pi].
    \end{dcases}
\end{equation}

Entonces calculamos la integral usando la definición,

\begin{align*}
    \int_{\gamma}^{}z^{2}\odif{z} & =
    \int_{0}^{3\pi}\gamma^{2}(t)\gamma^{\prime}(t)\odif{t},                                                      \\
    & = \int_{0}^{\pi}\pi^{2}\mathrm{e}^{2it}\odif{t} + \int_{\pi}^{3\pi}\bigl(t -
    2\pi\bigr)^{2}\odif{t},                                                                                      \\
    & = i\pi^{3} \int_{0}^{\pi}\mathrm{e}^{3it}\odif{t} +
    \eval{\dfrac{\bigl((t -
    2\pi)\bigr)^{3}}{3}}_{\pi}^{3\pi},                                                                           \\
    & = i\pi^{3}\Bigl(\int_{0}^{\pi}\cos\bigl(3t\bigr)\odif{t} + i
    \int_{0}^{\pi}\sin\bigl(3t\bigr)\odif{t}\Bigr) + \dfrac{2\pi^{3}}{3},                                        \\
    & = -\pi^{3}\Bigl(\eval{-\dfrac{1}{3}\cos\bigl(3t\bigr)}_{0}^{\pi}\Bigr) +
    \dfrac{2\pi^{3}}{3},                                                                                         \\
    & = -\dfrac{\pi^{3}}{3} - \dfrac{\pi^{3}}{3} + \dfrac{2\pi^{3}}{3},            \\
    & = 0.
\end{align*}

Ahora veamos si llegamos al mismo resultado usando el TFCC. Sea \(F(z) =
\dfrac{z^{3}}{3}\) la antiderivada, tal que

\begin{equation*}
    \int_{\gamma}^{}f(z)\odif{z} = F(\pi) - F(\pi) = 0.
\end{equation*}

\begin{proposition}
    Sea \(\Omega\footnote{Subconjunto abierto y conexo} \subseteq \mathbb{C}\) una
    región, \( f \colon \Omega \subseteq \mathbb{C} \to \mathbb{C}\) continua. Las
    siguientes proposiciones son equivalentes:

    \begin{enumerate}[label=(\roman*)]
        \item \(\exists F \colon \Omega \to \mathbb{C}\) analítica tal que

            \begin{equation*}
                F^{\prime}(z) = f(z),\qquad \forall z\in \mathbb{C}.
            \end{equation*}

        \item Si \(\gamma\) es una curva cerrada, i.e. \(\gamma(a) = \gamma(b)\), entonces

            \begin{equation*}
                \int_{\gamma}^{}f(z)\odif{z} = 0.
            \end{equation*}

        \item Si \(\gamma_{1} \colon [a, b] \to \Omega\) y \(\gamma_{2} \colon [c,
            d] \to \Omega\) curvas suaves por pedazos tales que \(\gamma_{1}(a) =
            \gamma_{2}(c)\) y \(\gamma_{1}(b) = \gamma_{2}(d)\), entonces

            \begin{equation*}
                \int_{\gamma_{1}}^{}f(z)\odif{z} = \int_{\gamma_{2}}^{}f(z)\odif{z},
            \end{equation*}

            i.e. la integral no depende de la curva.
    \end{enumerate}
\end{proposition}

\begin{proof}
    \vphantom{Texto para que no falle el ambiente de demostración chafa}
    \begin{enumerate}[label=""]
        \item (I) \(\;\implies\;\) (II) Supongamos que \(f\) tiene primitiva, entonces
            por el TFCC,

            \begin{align*}
                \int_{\gamma}^{}f(z)\odif{z} & = F(\gamma(b)) - F(\gamma(a)),     \\
                & = F(\gamma(b)) - F(\gamma(b)) = 0.
            \end{align*}
        \item (II) \(\;\implies\;\) (III) Notemos que \(\gamma_{1} - \gamma(2)\) es una curva
            cerrada y por (II),

            \begin{equation*}
                0 = \int_{\gamma_{1} - \gamma_{2}}^{}f(z)\odif{z} = \int_{\gamma_{1}}^{}f(z)\odif{z} -
                \int_{\gamma_{2}}^{}f(z)\odif{z}.
            \end{equation*}
            \classdate{2026-04-13}
        \item  (II) \(\;\implies\;\) (III) Sea \(z_{0} \in U\) fijo y sea otro \(z \in
            U\). Definimos

            \begin{equation}
                F(z) = \int_{z_{0}}^{z}f(z)\odif{z},
            \end{equation}

            donde por abuso de notación, la integral de la derecha denota la integral a lo largo de
            cualquier curva en \(U\) que una \(z_{0}\) con \(z\), pues por hipótesis la integral no
            depende de la curva.

            Sea \(r > 0\) tal que \(B_{r}(z_{0}) \subseteq U\), con \(U\) abierto y conexo. en
            particular, si tomamos \(h \in \mathbb{R}\) con \(0 < h < r\), se tiene que \(z_{0} + h,
            z_{0} + ih\in B_{r}(z_{0})\). Entonces,

            \begin{equation*}
                \pdv{F(z)}{x} = \lim_{h \to 0} \dfrac{F(z + h) - F(z)}{h} = \lim_{h\to
                0}\dfrac{1}{h}\Biggl(\int_{z_{0}}^{z + h}f(w)\odif{w} -
                \int_{z_{0}}^{z}f(w)\odif{w}\Biggr).
            \end{equation*}

            Donde la integral de \(z_{0}\) a \(z +  h\) es

            \begin{align*}
                \int_{z_{0}}^{z + h}f(w)\odif{w} &= \int_{z_{0}}^{z}f(w)\odif{w} + \int_{z}^{z +
                h}f(w)\odif{w},\nonumber\\
                \implies \int_{z}^{z + h}f(w)\odif{w} &= \int_{z_{0}}^{z + h}f(w)\odif{w} -
                \int_{z_{0}}^{z}f(w)\odif{w},
            \end{align*}

            tal que

            \begin{equation*}
                \pdv{F(z)}{x} = \lim_{h\to 0}\dfrac{1}{h}\int_{z}^{z + h}f(w)\odif{w} = \lim_{h\to
                0}\int_{0}^{d}f(z + t)\odif{t}.
            \end{equation*}

            Ahora vamos a tomar la parte real e imaginaria de \(f(z + t)\), de tal forma que las
            siguientes integrales son de funciones reales,

            \begin{equation*}
                \dfrac{1}{h}\int_{0}^{h}\Re f(z + t)\odif{t}\quad \wedge\quad
                \dfrac{1}{h}\int_{0}^{h}\Im f(z + t)\odif{t}.
            \end{equation*}

            Recordando el teorema del valor medio para integrales (T.V.M.I.). \(f \colon
            [a, b] \to \mathbb{R}\;\exists c\in[a, b]\) tal que

            \begin{equation*}
                f(c) = \Bigl(\int_{a}^{b}f(x)\odif{x}\Bigr)\bigl(b - a\bigr).
            \end{equation*}

            Entonces, \(\exists h^{\prime}, h^{\prime\prime}\in[0, h]\) tal que

            \begin{align*}
                \dfrac{1}{h} \int_{0}^{h}\Re f(z + t)\odif{t} &= \dfrac{1}{h}\Re f(z + h^{\prime})h =
                \Re f(z + h^{\prime}),\\
                \dfrac{1}{h} \int_{0}^{h}\Im f(z + t)\odif{t} &= \dfrac{1}{h}\Im f(z + h^{\prime\prime})h =
                \Im f(z + h^{\prime\prime}).
            \end{align*}

            Si \(h\to 0\), entonces \(h^{\prime}, h^{\prime\prime} \to 0\). Como es continua
            obtenemos que

            \begin{align*}
                \lim_{h^{\prime} \to 0} \Re f(z + h^{\prime}) = \Re f(z),\\
                \lim_{h^{\prime\prime} \to 0} \Im f(z + h^{\prime\prime}) = \Im f(z).
            \end{align*}

            Si \(h^{\prime\prime\prime} < \min(h^{\prime}, h^{\prime\prime})\),

            \begin{equation*}
                \pdv{F(z)}{x} = \lim_{h^{\prime\prime\prime}\to 0}\dfrac{F(z +
                h^{\prime\prime\prime}) - F(z)}{h^{\prime\prime\prime}} = \Re f(z) + i \Im f(z) =
                f(z).
            \end{equation*}

            De manera análoga,

            \begin{equation*}
                \pdv{F(z)}{y} = \lim_{h\to 0}\dfrac{F(z +
                ih) - F(z)}{h} = \lim_{h\to 0}\dfrac{1}{h}\int_{z}^{z + ih}f(w)\odif{w} = \lim_{h\to
                0}\dfrac{1}{h}\int_{0}^{h}f(z + it)i\odif{t}= i f(z).
            \end{equation*}

            Por lo tanto, notemos que se satisface

            \begin{equation*}
                \pdv{F}{x} = i \pdv{F}{y},
            \end{equation*}

            i.e. cumplen las ecuaciones de Cauchy-Riemann en su versión compleja.
            Además, como \(f\) es continua, por ser analítica, \(\pdv{F}{x}\) y
            \(\pdv{F}{y}\) son continuas, por lo que \(F\) es de clase \(C^{1}\), lo
            cual implica que \(F\) es analítica.

            Y además, \(F = u + iv\),

            \begin{equation*}
                f(z) = \pdv{F(z)}{x} = \pdv{u(z)}{x} + i \pdv{v(z)}{x} = F^{\prime}(z).
            \end{equation*}

            Por lo tanto, \(F^{\prime} = f\).
    \end{enumerate}
\end{proof}
\end{document}
